\section{Related Work}
\label{sec:related}

\subsection{ReCom and GerryChain}

The \recom{} (\emph{Recombination}) Markov chain was introduced by
\citet{deford2021recombination} as a method for sampling redistricting
plans with a provably well-defined stationary distribution.
At each step, the chain selects a pair of adjacent districts, merges
them into a single region, generates a random spanning tree of the merged
region using Wilson's algorithm, finds all balanced bipartitions of that
tree, and selects one uniformly at random.
The chain's stationary distribution is uniform over the space of
population-balanced redistricting plans, up to the spanning-tree bias
introduced by graph structure.

\gc{}~\citep{gerrychain2021} is the canonical open-source Python
implementation of \recom{}, developed and maintained by the MGGG
Redistricting Lab at Tufts University.
It has been used in expert reports submitted in redistricting litigation
in North Carolina~\citep{herschlag2020quantifying}, Wisconsin, and
Pennsylvania~\citep{chikina2017assessing}, and forms the computational
backbone of the ensemble-based auditing literature.
\gc{} is designed for correctness and research flexibility over raw
throughput; it runs at approximately 21--36 steps per second for
congressional-scale states.

\citet{autry2021multiscale} extended the \recom{} framework to a
multiscale forest variant, improving mixing properties for large $k$.
\citet{carter2019merge} analyzed the \emph{merge-split} chain, a closely
related Markov chain with distinct stationary distribution properties.

\subsection{Wilson's Spanning Tree Algorithm}

\citet{wilson1996generating} gave the first efficient algorithm for
generating uniform random spanning trees of an arbitrary graph.
Wilson's algorithm performs a loop-erased random walk from an arbitrary
vertex to the already-constructed tree, erasing loops as they form and
adding the loop-erased path to the tree.
The algorithm runs in expected time equal to the \emph{mean cover time}
of a simple random walk on the graph.

The loop-erasing property ensures that the resulting spanning tree is
drawn uniformly from the set of all spanning trees of the graph (by
Kirchhoff's matrix-tree theorem), making it a direct tool for \recom{}
proposals.
\citet{aldous1991random} proved that the cover time of a random walk on
a planar graph with $m$ vertices is $O(m\log m)$, giving Wilson's
algorithm $O(m\log m)$ expected time on census-tract adjacency subgraphs.

The loop-erased random walk has a long independent history
dating to~\citet{lawler1980loop}, who studied it as a model in
statistical mechanics.
\citet{propp1998exact} showed that Wilson's algorithm is a special case
of their general \emph{coupling from the past} framework.

\subsection{Rust for High-Performance Scientific Computing}

Rust's combination of zero-cost abstractions, safe concurrency, and
predictable memory layout has made it an increasingly attractive language
for computationally intensive scientific applications.
The Rayon data-parallelism library~\citep{rayon2024} provides
fork-join parallelism with work-stealing scheduling and a threading model
that maps cleanly onto independent MCMC chains.
\citet{matsakis2014rust} describe the language's ownership system, which
eliminates data races by construction---a critical property for parallel
MCMC where chains share no mutable state.

Within the \texttt{redist} ecosystem, the Rust codebase has achieved
approximately $213\times$ speedup over the Python pipeline for the core
METIS partitioning workflow~\citep{dellaLibera2026b1}.
The \re{} crate extends this philosophy to ensemble sampling.

\subsection{Portfolio Context}

The G-track papers in this portfolio provide the ensemble methodology and
comparison results that \re{} operationalizes.
\citet{dellaLibera2026g0} establishes the G-track framework for comparing
deterministic redistricting algorithms to \gc{} ensembles.
\citet{dellaLibera2026g1} provides direct percentile placement of the
\ar{} plans within \gc{} ensembles for six states.
\citet{dellaLibera2026g4} formalizes the $\rhat$, $\ess$, and Hamming
autocorrelation diagnostics implemented in the \texttt{redist} binary.
The H.1 paper~\citep{dellaLibera2026h1} establishes the bisection-ensemble
connection that motivates real-time ensemble integration.
\re{} closes the final methodological gap: G-track results previously
depended on the Python \gc{} installation; \re{} makes those results
reproducible from the \texttt{redist} Rust binary alone.
